\documentclass[12pt]{article}

\usepackage[american]{babel}
\usepackage[margin=1in]{geometry} 
\usepackage{amsmath,amsthm,amssymb}
\usepackage{multicol}
\setlength{\columnsep}{1cm}
\usepackage{mathrsfs}
\usepackage{tikz}
\usepackage{color}
\usepackage{lipsum} % for filler text
\usepackage{fancyhdr}
\usepackage{graphicx}
\pagestyle{fancy}
\usepackage[svgnames]{xcolor}
\usepackage[colorlinks=true, allcolors=DeepPink]{hyperref}
\usepackage{hyperref}
\usepackage{float}
\usepackage{subcaption}
\usepackage{xcolor}
\usepackage{wrapfig}

\newcommand{\N}{\mathbb{N}}
\newcommand{\Z}{\mathbb{Z}}
\newcommand{\ep}{\varepsilon}
\newcommand{\pa}{\partial}

\usepackage{titling}
\usepackage{hyperref}

\setlength{\droptitle}{-8em}   % This is your set screw
\setlength{\headheight}{16pt}

\begin{document}

\title{AMATH 732 Project Ideas}
\author{the royal you (you as in we)}
\maketitle

\section{Introduction}
i love latex
yay $\epsilon$ $\varepsilon$

\begin{table}[H]
    \centering
    \begin{tabular}{p{3in}|p{3in}}
         Stochastics (Individual Based Model)  & Squid Pop'n and Overfishing \\ \hline
         Could either do something related to a disease going through a pop'n of humble squid OR could do some kind of squid fishery thing ---- could do a disease spreading through a squid fishery with squeeze points and different disease dynamics -- something with changing the size of the space they can travel and also increasing the pop'n rapidly(ish) & Overfishing of other fish changes the squid carrying capacity - how does that change over time? \\
    \end{tabular}
    \caption{shoutout projects}
    % \label{tab:placeholder}
\end{table}

\textbf{Stochastics (Individual Based Model):}
\begin{table}[H]
    \centering
    \begin{tabular}{p{1.5in}|p{1.5in}|p{1.5in}|p{1.5in}}
         Resources  & Tasks & Attractiveness & Flaws \\ \hline
         literature, available code & modeling, code & link to stochastic modeling, accessible, would be cool to model swimming/the fluid motion of squid in a pop'n, cooler pictures & will require a concrete focus, lots of debugging \\
    \end{tabular}
    \caption{shoutout squid}
\end{table}


Squid Pop'n and Overfishing:
\begin{table}[H]
    \centering
    \begin{tabular}{p{1.5in}|p{1.5in}|p{1.5in}|p{1.5in}}
         Resources  & Tasks & Attractiveness & Flaws \\ \hline
         well studied i think, lots of resources  & modeling, code, analysis of model & Accessible and ODE-based (yay ODEs), much more already known analytical theory and graphing directions & will need some work to be novel, could get super complicated super quickly, mostly likely will spiral into a PDE model \\
    \end{tabular}
    \caption{shoutout fish}
    % \label{tab:placeholder}
\end{table}

potential ideas: population vs. temperature, population vs. populations of other fish?
\\$\Rightarrow$ I think we can combine these into one model! Would be super interesting (-Madi)
%also i fixed the tables #ImLearning
% i didnt know you could do arrows in latex


\section{POTENTIAL RESOURCES} %tysm !!!!
\begin{table}[H]
    \centering
    \begin{tabular}{p{2.0in}|p{1.0in}|p{2.5in}}
         Title  & Author & Notes \\ \hline
        \href{https://onlinelibrary.wiley.com/doi/abs/10.1111/j.1365-2486.2010.02223.x}{Modelling climate-change-induced nonlinear thresholds in cephalopod population dynamics} & Andre et al., 2010 & blah \\ \\
        \href{https://link.springer.com/article/10.1007/s11160-007-9077-3}{The potential impacts of climate change on inshore squid: biology, ecology and fisheries} & Pecl \& Jackson, 2007 & bleh\\ \\
        \href{https://www.annualreviews.org/content/journals/10.1146/annurev.ecolsys.36.102003.152644}{Individual-Based Modeling of Ecological and Evolutionary Processes} & Donald L. DeAngelis and Wolf M. Mooij, 2005 &  Review from `05 of individual-based models in ecology \\ \\

        \href{https://cdnsciencepub.com/doi/abs/10.1139/f93-015}{Growth Rate Variation and Larval Survival: Inferences from an Individual-Based Size-Dependent Predation Model} & Rice et al., 1993 & Use of a individual-based model to predict recruitment success of a cohort of fish based on length and growth rate\\ \\

        \href{https://www.degruyterbrill.com/document/doi/10.1515/9781400850624/html}{Individual-based Modeling and Ecology} & Volker Grimm and Steven F. Railsback, 2003 (book) & found a pdf online. book seems to be written for ecologists---mathematical detail seems a little sparse\\ \\ 
        \href{https://academic.oup.com/beheco/article/16/1/178/206370}{Density distribution and size sorting in fish schools: an individual-based model} & Charlotte K. Hemelrijk , Hanspeter Kunz, 2004 & fish schooling \\ \\ 

        \href{https://onlinelibrary.wiley.com/doi/abs/10.1046/j.1365-2419.1997.00037.x}{
An individual based model of fish school reactions: predicting antipredator behaviour as observed in nature} &RUNE VABØ, LEIF NØTTESTAD, 2003 & cool
    
    \end{tabular} 
    \caption{shoutout literature}
    % \label{tab:placeholder}
\end{table}

\section{TASKS (\& who's doing them)}
\begin{itemize}
    \item madi: build model
    \item cole:
    \begin{itemize}
        \item set up github / version control
        \item data analysis / statistics??
        \item making powerpoint fun
    \end{itemize}
    \item sarah: 
        \begin{itemize}
            \item writing?
            \item part of data analysis
            \item part of presentation
            \item is our model realistic
        \end{itemize}
        
\end{itemize}

\section{CORE MODEL}
\begin{itemize}
    \item FISH
    \begin{itemize}
        \item fish population / concentration
        \item how do they move (water current, brownian motion, etc.)
        \item amount of schooling / fish schooling mechanic
    \end{itemize}
    \item SQUID
    \begin{itemize}
        \item jet propulsion?
        \item choose fish based on multiple criteria? need mechanic
        \begin{itemize}
            \item how close to fish
            \item how close fish is to other fish
        \end{itemize}
    \end{itemize}
\end{itemize}

\subsection{GOALS}
\begin{itemize}
    \item how much does schooling matter
    \item how much does the wait time between eating matter
    \item how much does squid-fish coupling matter?
    \item fish population / concentration over time
\end{itemize}

\subsection{FUTURE WORK/SUPPLEMENTAL ACTIVITIES}
\begin{itemize}
    \item multiple squid / shoal of squid?
    \item water current / fish propulsion
\end{itemize}



\end{document}
